\chapter{The CnaTests Architecture}
\label{ch:cnatests-architecture}
\label{ch:verification-methodology}
\index{CnaTests}\index{GoogleTest}\index{CTest}\index{test architecture}

CNA's test system is not a directory of independent unit-test programs.  It is a composition
of one large GoogleTest binary, hundreds of standalone renderer executables, process-isolation
harnesses, negative compilation, and script-based structural gates.  Understanding which layer
produced a result is a prerequisite for interpreting that result.

\section{One discovered GoogleTest binary}

\texttt{cmake/UnitTests.cmake} recursively globs C\texttt{++} test sources from three roots:
non-renderer modules, renderer modules, and the repository-level \texttt{tests} directory.  The
result becomes exactly one executable, \texttt{CnaTests}.  There is no GoogleTest executable per
framework module.

The broad glob is narrowed by nine explicit \texttt{list(FILTER ... EXCLUDE REGEX)} operations.
They remove standalone Glide ABI programs and minimal-link probes, conditionally remove five
FFmpeg video tests, remove Net and GamerServices trees when networking is disabled, exclude
POSIX-process tests on Windows, Emscripten, and Android, and select only the configured
renderer-specific test directory.  One exclusion group even checks itself with a fatal error
if Net sources survived a Net-disabled configure.  Filtering is therefore part of the build
contract, not housekeeping.

GoogleTest cases are exposed to CTest with
\texttt{gtest\_discover\_tests(... DISCOVERY\_MODE PRE\_TEST)}.  Discovery runs the built binary
at test time, so the complete CTest list does not exist at configure time.  Its working directory
is pinned to the source root after cases failed to find \texttt{tests/assets} when CTest launched
them from the runtime-output directory.  Both details explain why a source-only count and a
\texttt{ctest -N} count answer different questions.

\section{One configure, one renderer set}

The ordinary configuration enters one chosen renderer family.  Alpha.1 also accepts a compatible
semicolon-separated \texttt{CNA\_GRAPHICS\_RENDERERS} set; the selected singular identity remains
the default.  Renderer-local source filtering tests membership in
\texttt{CNA\_RENDERER\_IDENTITIES}, so a multi-renderer \texttt{CnaTests} build retains the suites
for every compiled family rather than only the default.  Metal and Glide still contribute some
portable policy suites even when their native implementation is gated.  No supported tuple
therefore implies ``all renderers,'' and a result that omits both the compiled set and the active
runtime identity cannot be interpreted safely.

There are three layers of renderer gating:

\begin{enumerate}
  \item source filtering decides which compiled renderer families' GoogleTests enter
        \texttt{CnaTests};
  \item CMake conditions decide which standalone executables and CTests are registered;
  \item runtime pre-flight may return the project-wide skip code~77 when no display is usable.
\end{enumerate}

That last result is deliberately SKIPPED rather than FAILED.  It keeps headless machines from
turning absence of a display into a renderer defect, but it also means a green CTest summary may
contain no executed pixel comparison.  Reports must retain skip counts and the display setup.

\section{Standalone examples are a separate corpus}

Files named \texttt{*\_test.cpp} under renderer \texttt{examples} usually do not contain
GoogleTest at all.  Per-renderer CMake macros build them one executable at a time, each with its
own \texttt{main()}, then optionally register it with CTest.  Many consume the same public
fixture from \texttt{modules/graphics/examples}, allowing one renderer to expose a defect while
the others serve as controls.

Building and registering are independent acts.  Some renderers build far more executables than
they register; others add Python audit registrations that are not executables.  Labels are also
uneven: \texttt{GraphicsSmoke} is cross-renderer, while several large renderer families label
only one or two of hundreds of registrations.  A command such as \texttt{ctest -L EasyGL}
therefore cannot be assumed to select the EasyGL corpus.

\section{Why some tests need another process}

Several failures concern process-global state: SDL shutdown, the audio mixer's global owner,
GamerServices dispatch, and network peers.  CNA builds small helper executables for these cases
and embeds their absolute paths into \texttt{CnaTests}.  A GoogleTest case spawns the helper and
judges its result, gaining a clean address space and making crashes or shutdown ordering
observable without poisoning later cases.

Other helpers are diagnostic tools rather than isolation tests: an XNB audio metadata dumper,
XWB inspector, FNA-reference dump, and Devices microbenchmark.  Their presence beside harnesses
does not make them part of every test run.  The governing question is always who invokes the
target and under what configuration.

\section{Compilation failure can be the oracle}

\texttt{cna\_strict\_xna\_api\_leak\_check} is excluded from the normal build.  Its CTest command
invokes the build system for that target and is marked \texttt{WILL\_FAIL TRUE}.  The test passes
only when the compiler rejects the forbidden API use.  This inverted polarity turns a namespace
and surface rule into executable evidence: successful compilation would be the regression.

Fourteen module probes apply the complementary idea to linking.  Each builds a minimal consumer
of one module alias, then a link-closure test inspects the generated \texttt{link.txt} against a
module-specific forbidden-library expression.  The probe asks whether the public target works;
the closure gate asks whether it works for the right architectural reason.

\section{Discovery, selection, execution, verdict}

A reliable test report separates four stages:

\begin{description}
  \item[Discovery] Which sources and runtime-expanded GoogleTest cases were present?
  \item[Selection] Which filter, label, renderer, and feature options chose the cases?
  \item[Execution] Did the intended executable, display, GPU, browser, or compatibility layer
        actually engage?
  \item[Verdict] Which oracle declared pass, fail, or skip?
\end{description}

This model explains most apparent contradictions in CNA's historical test prose.  A source can
exist but be filtered out; a target can build but not be registered; a CTest can register but
skip; a program can pass while running through the wrong graphics translation layer.  The rest
of this Part examines the numbers and oracles that make those stages auditable.
